% -------------------------------------------------------------------------
% WBS ID: RES-DAT-SRC
% Deliverable: Earthquake and Tsunami Source Data
% Status: In progress
% Evidence status: Regional model source-data requirements established; final dataset pending
% Review status: Not reviewed
% Last updated: 2026-07-19
% Dependencies: RES-PHY-SRC, RES-MOD-SRC, RES-MOD-SPEC
% Downstream interfaces: RES-NUM-SRC, RES-VER-SPEC, Software source ingestion
% -------------------------------------------------------------------------

\section{RES-DAT-SRC --- Earthquake and Tsunami Source Data}
\label{sec:res-dat-src}

\subsection{Purpose and Scope}
\label{sec:res-dat-src-purpose}

This Deliverable defines the earthquake event metadata, finite-fault
source candidates and bed-deformation data required to reconstruct the
2011 Tōhoku tsunami in the regional 2D NLSWE model. It covers source
data selection and projection into model-ready bed increments. It does
not choose a structural damage model, a local barrier geometry or a
solver implementation.

\subsection{Research Questions}
\label{sec:res-dat-src-research-questions}

The source-data decision asks:

\begin{enumerate}
  \item Which event metadata define the authoritative epicentre,
    origin time, magnitude and depth for corridor construction?
  \item Which finite-fault or tsunami-source reconstructions are
    candidates for regional bed forcing?
  \item Which observations were used to infer a source and therefore
    cannot later be claimed as independent validation evidence?
  \item Which data fields must Software receive to generate
    $\Delta b_m(x,y)$ and rupture-time functions reproducibly?
\end{enumerate}

\subsection{Terminology and Definitions}
\label{sec:res-dat-src-definitions}

% TODO: Define the technical terminology, symbols, quantities and
% classifications used in this Deliverable.

\subsection{Literature Evidence}
\label{sec:res-dat-src-literature-evidence}

% TODO: Record evidence from peer-reviewed papers, authoritative datasets,
% standards, technical reports and primary documentation.
%
% Do not create a detached literature-summary list. Explain how each source
% supports, contradicts or limits a project decision.

\textcite{usgs2011tohokuEvent} is retained as the event-metadata
source for the 2011 Great Tōhoku earthquake: origin time
2011-03-11 05:46:24 UTC, epicentral coordinates near
$38.297^{\circ}\mathrm{N}$, $142.373^{\circ}\mathrm{E}$, depth
$29.0~\mathrm{km}$ and moment magnitude $M_{\mathrm{w}}=9.1$. Source
role: supports event identity, epicentre coordinate and corridor
origin. Project conclusion: these metadata define
$P_{\mathrm{e}}=(\lambda_{\mathrm{e}},\phi_{\mathrm{e}})$ until a
GIS-approved event-source point is selected. Limitation: event
metadata alone do not define the spatial seabed displacement field.

\textcite{usgsFiniteFaults} and \textcite{hayes2011rapid} are retained
for the USGS finite-fault source-data family and rapid Tōhoku source
characterisation. Source role: supports candidate finite-fault model
fields: subfault geometry, slip, rake, rupture time and rise time.
Project conclusion: the USGS finite-fault family is a baseline source
candidate for regional moving-bed forcing. Limitation: the finite-fault
model must still be converted into hydrodynamic bed increments and
tested against observation sets not used in the inversion.

\textcite{fujii2011tsunamisource} and \textcite{wei2013japan} remain
retained source-reconstruction references. Source role: support and
limit the tsunami-source calibration choices used in Tōhoku hindcasts.
Project conclusion: these sources are comparison and provenance
anchors when deciding whether the USGS finite-fault source or an
alternative tsunami-inversion source better supports the selected
Kamaishi and Sendai corridors. Limitation: any observation record used
to tune a source shall be marked $\mathcal{D}_{\mathrm{cal}}$ in
`RES-DAT-OBS`.

\subsection{Governing Relations and Quantitative Definitions}
\label{sec:res-dat-src-governing-relations}

% TODO: Record relevant equations, dimensionless groups, empirical models,
% extraction rules or quantitative definitions.
%
% Use align environments for displayed equations.
% Every displayed equation must have a unique \label{eq:...}.
% Do not place punctuation after displayed equations.

\subsection{Required Data Variables and Units}
\label{sec:res-dat-src-required-data-variables-units}

The dynamic moving-bed regional source requires, for each source
component or subfault $m$:

\begin{itemize}
  \item geometry, location, depth, length and width;
  \item strike, dip and rake;
  \item slip or displacement amplitude;
  \item rupture timing and rise-time function $R_m(t)$;
  \item projected bed increment $\Delta b_m(x,y)$;
  \item provenance and uncertainty metadata.
\end{itemize}

\subsection{Candidate Data Sources}
\label{sec:res-dat-src-candidate-data-sources}

The current source candidates are:

\begin{itemize}
  \item USGS event metadata for origin time, magnitude, depth and
    epicentre \parencite{usgs2011tohokuEvent};
  \item USGS finite-fault products and rapid finite-fault source
    characterisation \parencite{usgsFiniteFaults,hayes2011rapid};
  \item Fujii--Satake source reconstruction
    \parencite{fujii2011tsunamisource};
  \item DART-constrained and nested-model source workflows described
    by \textcite{wei2013japan}.
\end{itemize}

The source-data table in `RES-DAT-CAS` records each candidate's role,
downstream deliverable and limitation.

\subsection{Spatial and Temporal Coverage}
\label{sec:res-dat-src-spatial-temporal-coverage}

Spatial coverage shall include the full rupture footprint and enough
surrounding seafloor for interpolation of dynamic bed increments onto
the regional mesh. Temporal coverage shall include the origin time,
subfault rupture onset and rise-time function for every source
component retained in the moving-bed source.

\subsection{Coordinate and Datum Requirements}
\label{sec:res-dat-src-coordinate-datum-requirements}

Raw source coordinates shall be retained in longitude/latitude and
source-native depth convention. The solver-facing product shall be
projected into the same horizontal metric frame as `RES-DAT-GEO` and
shall use positive-upward bed increments compatible with
\cref{eq:res-dat-src-stage-bed-update}. Any conversion from fault slip
to seabed displacement shall record the elastic model, coordinate
frame, vertical sign convention and interpolation method.

\subsection{Data Processing and Transformation}
\label{sec:res-dat-src-data-processing-transformation}

The hydrodynamic source data shall be delivered as bed increments
compatible with:

\begin{align}
  b_i(t_s)
  &=
  b_{0,i}
  +
  \sum_m
  \Delta b_{m,i}
  R_m(t_s)
  \label{eq:res-dat-src-stage-bed-update}
\end{align}

where $t_s$ is the relevant Runge--Kutta stage time. Static source
initialisation may be retained as a verification or fallback case, but
the regional model specification treats dynamic moving bathymetry as the
selected mathematical source representation.

\subsection{Uncertainty, Provenance and Licensing}
\label{sec:res-dat-src-uncertainty-provenance-licensing}

The source manifest shall identify whether each observation used to
create or tune a source is calibration, validation or diagnostic
evidence. Finite-fault uncertainty shall be retained at least as source
version, inversion method, data types used, subfault discretisation,
rupture timing assumptions and published update date where available.
Licence and access notes shall be recorded before any finite-fault data
are stored in a project dataset bundle.

\subsection{Data Acceptance Criteria}
\label{sec:res-dat-src-data-acceptance-criteria}

A source dataset is accepted for regional replay only when it supplies
or can defensibly derive:

\begin{itemize}
  \item event origin metadata and source version;
  \item subfault geometry, strike, dip, rake, slip and depth;
  \item rupture timing or a documented static-source fallback;
  \item projected bed increments $\Delta b_{m,i}$ on the regional
    mesh;
  \item calibration/validation-use classification for all observation
    records used in the source construction;
  \item provenance, licence and transformation metadata.
\end{itemize}

\subsection{Candidate Approaches}
\label{sec:res-dat-src-candidate-approaches}

% TODO: Define the credible candidate approaches considered.

\subsection{Critical Comparison}
\label{sec:res-dat-src-critical-comparison}

% TODO: Compare candidates using physically and project-relevant criteria.
% Do not select an approach solely on popularity or implementation ease.

\subsection{Assumptions and Applicability}
\label{sec:res-dat-src-assumptions}

% TODO: State assumptions and the conditions under which they are valid.

\subsection{Limitations and Failure Conditions}
\label{sec:res-dat-src-limitations}

% TODO: State where the evidence, model or selected approach ceases to be
% reliable or applicable.

\subsection{Project-Specific Conclusions}
\label{sec:res-dat-src-conclusions}

The Tōhoku source decision is source-data finalisation, not solver
implementation. The USGS event metadata define the current corridor
origin and event identity. USGS finite-fault, Hayes rapid source,
Fujii--Satake and DART-constrained source workflows remain source
candidates until their projected bed increments are compared against
the selected offshore, nearshore and terrestrial observation split.

\subsection{Selected Approach and Rationale}
\label{sec:res-dat-src-selected-approach}

Use a source-candidate manifest with USGS event metadata as the event
identity baseline and USGS finite-fault/Fujii/DART-constrained
reconstructions as candidate bed-forcing products. The selected
hydrodynamic source shall be the candidate that best supports the
Kamaishi and Sendai corridor validation split without reusing
calibration records as independent validation.

\subsection{Unresolved Decisions}
\label{sec:res-dat-src-unresolved-decisions}

% TODO: Record unresolved questions, required evidence and decision owners.

The source-data variables required by the mathematical model are
established. The exact final fault-source dataset, any comparison
against independent source reconstructions, final rupture-time
functions and the accepted source-to-bed projection remain open
data-selection tasks.

\subsection{Requirements and Handoffs}
\label{sec:res-dat-src-requirements}

% TODO: State requirements passed to other Research Domains, Software,
% verification activities or later execution work.

The selected source dataset shall be marked as calibration,
validation-support or diagnostic evidence so `RES-VER` can distinguish
source-fit records from independent validation records.

Software shall start a source-data ingestion path in parallel with the
remaining research: source manifest, fault-model parser, coordinate
projection, bed-increment rasterisation, structured source-output
schema and reproducible moving-bed replay fields for the regional
solver.

\subsection{Deliverable Completion Record}
\label{sec:res-dat-src-completion-record}

% TODO: Record whether the Deliverable is:
%
% - Not started
% - In progress
% - Draft complete
% - Reviewed
% - Accepted
%
% Acceptance requires:
%
% - the research questions to be answered;
% - claims to be supported by appropriate evidence;
% - assumptions and limitations to be explicit;
% - the project-specific conclusion to be stated;
% - downstream requirements to be traceable;
% - unresolved decisions to be closed or formally deferred.
